\section{Auslander regular ring}
Let $A$ be a noetherian ring (not necessarily commutative), and $\mathrm{Mod}_{f}(A)$, $\mathrm{Mod}_{f}^{op}(A)$ the abelian category of left/right finitely generated $A$-module.

For $M \in \mathrm{Mod}_{f}(A)$, the $\textbf{projective dimension}$ of $M$ is
$$\mathrm{pd}(M) = \text{minimum of length of projective resolution.}$$

We say $A$ is $\textbf{of finite global dimension}$, if $\exists w \in \mathbb{Z}_{>0}$ s.t for all $M \in \mathrm{Mod}_{f}(A)$, have $\mathrm{pd}(M) \le w$.

We now assume that $A$ is finite type over $k$ and of finite global dimension. Then $\mathrm{RHom}(M,A) \in \mathbf{D}^{b}(\mathrm{Mod}^{op}_{f}(A))$, the derived category of bounded complexes in $\mathrm{Mod}_{f}^{op}(A)$.

\begin{definition}
    \quad
    \begin{enumerate}
        \item The $\textbf{grade number}$ $j(M) := \min\left\{k \middle| \mathrm{Ext}^{k}(M,A) \neq 0\right\}$, we sometimes denote it as  $j_{A}(M)$ as well. In particular, define $j(M) = +\infty$ if $M = 0$.
        \item For a noetherian and of finite global dimension ring $A$, we say $A$ is $\textbf{Auslander regular}$ (AR for short) if for all $M \in \mathrm{Mod}_{f}(A)$ and $v \in \mathbb{Z}_{\ge 0}$, given any $N \subseteq \mathrm{Ext}^{v}(M,A)$ submodule, have $j(N) \ge v$.
    \end{enumerate}
\end{definition}

Let $S$ be a commutative $k$-algebra of finite type, we say $S$ is $\textbf{regular}$ if $\Omega_{S/k}$ is locally free, see definition \cite[\href{https://stacks.math.columbia.edu/tag/00OD}{Tag 00OD}]{stacks-project}.

\begin{theorem}[Auslander]\label{theorem 7.2}
    Let $S$ be regular of dimension $m$, $M \in \mathrm{Mod}_{f}(S)$, then
    \begin{enumerate}
        \item $S$ is of finite global dimension;
        \item $j(M) + \dim(\mathrm{supp}(M)) = m$;
        \item $\dim(\mathrm{supp}(\mathrm{Ext}^{i}(M,S))) \le m - i$ for all $i$;
        \item $\dim(\mathrm{supp}(\mathrm{Ext}^{j(M)}(M,S))) = \dim(\mathrm{supp}(M))$.
    \end{enumerate}
\end{theorem}
\begin{proof}
    \mbox{}

    For (1), see \cite[Thm.~19.12.]{eisenbud1995commutative} or \cite[\href{https://stacks.math.columbia.edu/tag/00OC}{Tag 00OC}]{stacks-project} for local ring case.
    Then use \cite[\href{https://stacks.math.columbia.edu/tag/00O9}{Tag 00O9}]{stacks-project} to get global case, which is \cite[\href{https://stacks.math.columbia.edu/tag/00OE}{Tag 00OE}]{stacks-project}.

    For (2), (3) and (4), in the case of $m$ is even, see \cite[Prop.~6.4.]{schnell_dmodules_notes}.

    For general $m$, one can see \cite[Thm.~D.4.4.]{hotta2008d}, but this book refers other books for the proof of (2), (3) and (4).
\end{proof}
\begin{corollary}
    Regular commutative $k$-algebra of finite type are Auslander regular.
\end{corollary}

Let $X = \mathrm{Spec } R$ be a smooth-affine variety over $\mathbb{C}$, our next goal is to prove $\mathcal{D}_{X}$ is AR.

\section{Spectral sequences of filtered complex}

Let $F^{\bullet}$ be a bounded complex of $\mathbb{Z}$-modules:$$\cdots \to F^{i-1} \xrightarrow{\mathrm{d}} F^i \xrightarrow{\mathrm{d}} F^{i+1} \to \cdots$$with an increasing $\mathbb{Z}$-filtration: $F_j^i \subseteq F_{j+1}^i \subseteq F^i \quad (\forall j,i \in \mathbb{Z})$ satisfying:
\begin{enumerate}
    \item Exhaustive condition $\bigcup\limits_j F_j^i = F^i \quad \forall i$ ;
    \item Hausdorff condition $\bigcap_j F_j^i = 0 \quad \forall i$
    \item  Differential preserves filtration $\mathrm{d}(F_j^i) \subseteq F_j^{i+1} \, (\forall i,j)$ , which implies that $\dots\rightarrow F_{j}^{i-1}\xrightarrow{\mathrm{d}} F_{j}^{i}\xrightarrow{\mathrm{d}} F_{j}^{i+1}\rightarrow\dots$ is a complex.
\end{enumerate}
Then, to every triple $(k, i, j)$ we set:
\begin{itemize}
    \item $Z_j^i(k) = \{u \in F_j^i \mid du \in F^{i+1}_{j-k}\}.$
    \item $B^{i}_{j}(k)=F^{i}_{j}\cap \mathrm{d}(F^{i-1}_{j+k-1})=\mathrm{d}(Z^{i-1}_{j+k-1}(k-1))$
    \item $E_j^i(k) = \dfrac{Z_j^i(k) + F_{j-1}^i}{B_j^i(k) + F_{j-1}^i}$, $E_{k}^{-j,i+j}:=E_{j}^{i}(k)$, and $E_{\bullet}^i(k) = \bigoplus\limits_{j} E_j^i(k)$.
\end{itemize}

\begin{remark}
    Since we consider the increasing filtration $\dots\subseteq F_{j}\subseteq F_{j+1}\subseteq \dots \subseteq F$, while usually we consider a decreasing filtration for cochain complex, we can actually view $F_{j}$ as $F^{-j}$. That's the reason why we define $E_{k}^{-j,i+j}:=E_{j}^{i}(k)$.
\end{remark}
So $\mathrm{d}(Z_j^i(k)) \subseteq Z_{j-k}^{i+1}(k)$, which induces $$ E_k^{-j,i+j} := E_j^i(k) \xrightarrow{\mathrm{d}_{k}} E_k^{-j+k,i+j-k+1} = E_{j-k}^{i+1}(k) .$$Further, we have $$\dots\xrightarrow{\mathrm{d}_{k}}E_{j+k}^{i-1}(k)\xrightarrow{\mathrm{d}_{k}}E_{j}^{i}(k)\xrightarrow{\mathrm{d}_{k}}E_{j-k}^{i+1}(k)\xrightarrow{\mathrm{d}_{k}} \dots$$Consider $E^{i}_{\bullet}(k)=\bigoplus\limits_{j}E^{i}_{j}(k)$ and know that $\mathrm{d}_{k}:E^{i}_{\bullet}(k)\rightarrow E^{i+1}_{\bullet}(k)$, this makes $E_\bullet(k)$ a complex of graded $\mathbb{Z}$-modules.  And $E_{j}^{i}(0) = F^{i}_{j} /F_{j-1}^{i}$, i.e. $E^{i}_{\bullet}(0)=\mathrm{gr }_{\bullet}  F^{i}$ the associated graded complex of  $F^{i}$ for all $i$.

\medskip

$\{E_k^{-\ast,\bullet+\ast} = E_\ast^\bullet(k)\}_k$ is called the $\textbf{spectral sequence}$ of the filtered complex $F^{\bullet}$.

\medskip

\begin{lemma}\label{lemma 7.4}
    $E_\bullet^i(k+1) = H^i(E_{\bullet}^{\bullet}(k))$, more explicitly: $$E^{i}_{j}(k+1)\simeq \frac{\ker(E_{j}^{i}(k)\xrightarrow{\mathrm{d}_{k}}E_{j-k}^{i+1}(k) )}{\mathrm{Im}(E_{j+k}^{i-1}(k)\xrightarrow{\mathrm{d}_{k}}E_{j}^{i}(k))}.$$
\end{lemma}
\begin{proof}
    \quad

    (After verifying that $\mathrm{d}_{k}$ is well-defined) Consider $k=0$, we obtain that $E^{i}_{j}(0)= F^{i}_{j} /F^{i}_{j-1}$, so
    $$\mathrm{ker}(E^{i}_{j }(0)\rightarrow E^{i+1}_{j}(0))=\ker(F^{i}_{j} /F^{i}_{j-1}\rightarrow F^{i+1}_{j} /F^{i+1}_{j-1})=\left(\{u\in F^{i}_{j}:du\in F^{i+1}_{j-1}\}+F^{i}_{j-1}\right) /F^{i}_{j-1},$$
    and
    $$\mathrm{Im}(E^{i-1}_{j}(0)\rightarrow E^{i}_{j}(0))=\mathrm{Im}(F^{i-1}_{j} /F^{i-1}_{j-1}\rightarrow F^{i}_{j} /F^{i}_{j-1})=\mathrm{d}(F^{i-1}_{j})+F^{i}_{j-1} /F^{i}_{j-1},$$
    then we get $$E^{i}_{j}(1)\simeq \frac{\mathrm{ker}(E^{i}_{j }(0)\rightarrow E^{i+1}_{j}(0))}{\mathrm{Im}(E^{i-1}_{j}(0)\rightarrow E^{i}_{j}(0))}.$$

    For $k\ge 1$, we obtain that $F_{j-1}^{i}\subseteq F_{j+k-1}^{i}$$$E^{i}_{j}(k+1)=\frac{\{u\in F^{i}_{j}:du\in F^{i+1}_{j-k-1}\}+F^{i}_{j-1}}{F^{i}_{j}\cap \mathrm{d}(F^{i-1}_{j+k})+F^{i}_{j-1}}.$$

    For $\left(Z^{i}_{j}(k)+F^{i}_{j-1}\right) /F^{i}_{j-1}\xrightarrow{\mathrm{d}_{k}}\left(Z_{j-k}^{i+1}(k)+F^{i+1}_{j-k-1}\right) /F^{i+1}_{j-k-1}$,
    $$\ker \mathrm{d}_{k}:=\{u\in F^{i}_{j} /F_{j-1}^{i}: \exists x\in u\ s.t.\, \mathrm{d}x\in F^{i+1}_{j-k-1}\}=\left(Z^{i}_{j}(k+1)+F^{i}_{j-1}\right) /F^{i}_{j -1}$$
    And
    $$\mathrm{Im}\left(\left(Z^{i-1}_{j+k}(k)+F^{i-1}_{j+k-1}\right) /F^{i-1}_{j+k-1}\xrightarrow{\mathrm{d}_{k}}\left(Z^{i}_{j}(k)+F^{i}_{j-1}(k)\right) /F^{i}_{j-1}(k)\right)=\left(\mathrm{d}(Z^{i-1}_{j+k}(k))+F^{i}_{j-1}(k)\right) /F^{i}_{j-1}(k).$$
\end{proof}

\quad

\begin{definition}[Induced filtration on Cohomology]\label{def:induced_filtration_cohomology}
    The filtration of $F^{\bullet}$ induces a filtration of $H^i(F^{\bullet})$ by $$F_j H^i(F^\bullet) = \frac{Z_j^{i}(\infty) + \mathrm{d}F^{i-1}}{\mathrm{d}F^{i-1}}.$$where $Z_j^{i}(\infty) := \mathrm{Ker}(d) \cap F_j^i$.
\end{definition}
\begin{remark}
    For standard defintion, see \cite[Section~5.5]{Weibel1994HA}, or see Remark~\ref{rmk:induced-filtration-on-Ext} we proved two definition are equivalent there.
\end{remark}

In fact, we often view $Z^{i}_{j}(k)$ and $B^{i}_{j}(k)$ in $F^{i}_{j} /F^{i}_{j-1}$. Namely, we define
$$\eta^{i}_{j}:F^{i}_{j}\mapsto F^{i}_{j} /F^{i}_{j-1},$$
and
$$\overline{Z}^{i}_{j}(k):=\eta^{i}_{j}(Z^{i}_{j}(k)),$$
$$\overline{B}^{i}_{j}(k):=\eta^{i}_{j}(B^{i}_{j}(k)),$$
we obtain that $E^{i}_{j}(k)\simeq \frac{\overline{Z}^{i}_{j}(k)}{\overline{B}^{i}_{j}(k)}$. Moreover, we have a tower for all $i,j:$ $$
    0 \subseteq \overline{B}_j^i(0) \subseteq \overline{B}_j^i(1) \subseteq \cdots \subseteq \overline{B}_j^i(r) \subseteq \cdots \subseteq \overline{B}_j^i(\infty)
    \subseteq \overline{Z}_j^i(\infty) \subseteq \cdots \subseteq \overline{Z}_j^i(r) \subseteq \overline{Z}_j^i(r-1) \subseteq \cdots \subseteq \overline{Z}_j^i(0)
$$
We define the limit terms:
\[
    \overline{B}_j^i(\infty) := \varinjlim_{r} \overline{B}_j^i(r) \quad \text{and} \quad \overline{Z}_j^i(\infty) := \varprojlim_{r} \overline{Z}_j^i(r).
\]

\begin{lemma}
    Assuming the filtration is exhaustive and Hausdorff, we have:
    \[
        \overline{B}_j^i(\infty) = \eta_j^i(F_j^i \cap \mathrm{d}(F^{i-1}))
        \quad \text{and} \quad \overline{Z}_j^i(\infty)
        \, \textcolor{red}{\supseteq}\, \eta_j^i(Z_j^i(\infty)).
    \]
\end{lemma}
\begin{proof}
    For the first equality, recall that $\overline{B}_j^i(r) = \eta_j^i(F_j^i \cap \mathrm{d}(F_{j+r-1}^{i-1}))$.
    Since the filtration is exhaustive, we have $\bigcup_{r} F_{j+r-1}^{i-1} = F^{i-1}$.
    Thus $\bigcup_{r} \mathrm{d}(F_{j+r-1}^{i-1}) = \mathrm{d}(F^{i-1})$.
    Therefore, $\varinjlim_{r} \overline{B}_j^i(r) = \eta_j^i(F_j^i \cap \mathrm{d}(F^{i-1}))$.

    For the second formula, Hausdorff condition is NOT sufficient to ensure:
    \[
        \bigcap_{r}\eta_j^i(\{u \in F_j^i \mid \mathrm{d}u \in F_{j-r}^{i+1}\}) = \eta_j^i\left(\bigcap_{r}\{u \in F_j^i \mid \mathrm{d}u \in F_{j-r}^{i+1}\}\right)
    \]
    So we can only get an inclusion, not an equality.
\end{proof}
\begin{remark}
    See \cite[the discussion above Lemma~5.5.7.]{Weibel1994HA}.
\end{remark}

\begin{proposition}\label{prop:Boardman's_convergence}
    $\mathrm{gr}_{j}^{F}H^{i}(F)\simeq \frac{\eta_j^i(Z_j^i(\infty))}{\overline{B}^{i}_{j}(\infty)}$.
\end{proposition}

\begin{proof}
    See \cite[Lemma~5.5.7.(3)]{Weibel1994HA}, where $e^\infty_{pq}$ is just $\frac{\eta_p^q(Z_p^q(\infty))}{\overline{B}^{p}_{q}(\infty)}$ \cite[Lemma~5.5.7.(1)]{Weibel1994HA}, or see the proof below:

    Since we have defined induced filtration on cohomology, then
    \begin{align*}
        \mathrm{gr}_{j}^{F}H^{i}(F) & \simeq  F^j H^i(F^\bullet) \big/ F^{j-1} H^i(F^\bullet)                                                                                                                  \\
                                    & \simeq \left(\frac{Z_j^{i}(\infty) + \mathrm{d}F^{i-1}}{\mathrm{d}F^{i-1}} \right) \big/ \left(\frac{Z_{j-1}^{i}(\infty) + \mathrm{d}F^{i-1}}{\mathrm{d}F^{i-1}} \right) \\
                                    & \simeq \frac{Z_j^{i}(\infty) + \mathrm{d}F^{i-1}}{Z_{j-1}^{i}(\infty) + \mathrm{d}F^{i-1}}                                                                               \\
                                    & \simeq \frac{Z_j^{i}(\infty)}{Z_j^{i}(\infty) \cap (Z_{j-1}^{i}(\infty) + \mathrm{d}F^{i-1})}                                                                            \\
                                    & \simeq \frac{Z_j^{i}(\infty)}{Z_{j-1}^{i}(\infty) + (Z_j^{i}(\infty) \cap \mathrm{d}F^{i-1})}                                                                            \\
                                    & \simeq \frac{Z_j^{i}(\infty)}{Z_{j-1}^{i}(\infty) + (F_j^{i} \cap \mathrm{d}F^{i-1})}.
    \end{align*}
    The last three steps are obtained as follows:
    \begin{itemize}
        \item The third isomorphism uses the Second Isomorphism Theorem $(A+B)/B \cong A/(A \cap B)$ by setting $A=Z_j^i(\infty)$ and $B=Z_{j-1}^i(\infty) + \mathrm{d}F^{i-1}$.
        \item The fourth isomorphism uses the Modular Law $A \cap (B+C) = (A \cap B) + C$ since $C=Z_{j-1}^i(\infty) \subseteq Z_j^i(\infty)=A$.
        \item The last isomorphism holds because $Z_j^i(\infty) \cap \mathrm{d}F^{i-1} = (\Ker(\mathrm{d}) \cap F_j^i) \cap \mathrm{d}F^{i-1} = F_j^i \cap (\Ker(\mathrm{d}) \cap \mathrm{d}F^{i-1}) = F_j^i \cap \mathrm{d}F^{i-1}$ noting that $\mathrm{Im}(\mathrm{d}) \subseteq \Ker(\mathrm{d})$.
    \end{itemize}
    On the other hand,
    \begin{align*}
        \frac{\eta_j^i(Z_j^i(\infty))}{\overline{B}^{i}_{j}(\infty)} & \simeq \frac{Z_j^i(\infty) + F_{j-1}^i}{F_j^i \cap \mathrm{d}F^{i-1} + F_{j-1}^i}          \\
                                                                     & \simeq \frac{Z_j^i(\infty)}{Z_j^i(\infty) \cap (F_j^i \cap \mathrm{d}F^{i-1} + F_{j-1}^i)} \\
                                                                     & \simeq \frac{Z_j^i(\infty)}{F_j^i \cap \mathrm{d}F^{i-1} + (Z_j^i(\infty) \cap F_{j-1}^i)} \\
                                                                     & \simeq \frac{Z_j^i(\infty)}{F_j^i \cap \mathrm{d}F^{i-1} + Z_{j-1}^i(\infty)}.
    \end{align*}
    Thus $\mathrm{gr}_{j}^{F}H^{i}(F) \simeq \frac{\eta_j^i(Z_j^i(\infty))}{\overline{B}^{i}_{j}(\infty)}$.
\end{proof}

\quad

\begin{definition}\label{def:convergent_spectral_sequence}
    We say the spectral sequence $\{E_{\bullet}^{\bullet}(k)\}_k$ $\textbf{converges to}$ $H^i(F^{\bullet})$ if for all $i,j$, there exists $k = k(i,j)$ s.t. $\forall k' \ge k$, $$E_j^i(k') \cong \mathrm{gr}_j^F H^i(F^{\bullet}).$$
\end{definition}

\begin{definition}$\textbf{The regularity condition}:$
    if there exists a $\omega \in \mathbb{Z}_{\ge0}$ s.t.$$F_j^i \cap \mathrm{d}(F^{i-1}_{}) \subset \mathrm{d}(F_{j+\omega}^{i-1})$$for every pair $(i,j)$, then we say the filtered complex is $\omega$-regular.
\end{definition}

\begin{lemma}\label{lemma 7.8}
    If filtered complex $F^{\bullet}$ is $\omega$-regular and convergent in \cite[Convergence~5.2.11]{Weibel1994HA}'s sense, i.e. $\mathrm{gr}_{j}^{F}H^{i}(F) \simeq \frac{\overline{Z}_j^i(\infty)}{\overline{B}^{i}_{j}(\infty)}$. Then $$\mathrm{gr}^F_{\bullet} H^{i}(F^\bullet) \simeq E_{\bullet}^{i}(\omega+1), \quad \forall i$$In particular, the spectral sequence converges.
\end{lemma}

\begin{proof}
    First, we show that $E^{i}_{j}(\omega+k)\simeq E^{i}_{j}(\omega+1)$ for all $k\ge 1$.
    It suffices to show that the differentials $\mathrm{d}_r: E_{j+r}^{i-1}(r) \to E_j^i(r)$ are zero for all $r \ge \omega+1$.

    Any element in the image of $\mathrm{d}_r$ is represented by some $v \in F_j^i \cap \mathrm{d}(F_{j+r}^{i-1})$.
    By the regularity condition, we have
    $$ v \in F_j^i \cap \mathrm{d}(F^{i-1}) \subseteq \mathrm{d}(F_{j+\omega}^{i-1}). $$
    Since $r \ge \omega+1$, we have $j+\omega \le j+r-1$, so $F_{j+\omega}^{i-1} \subseteq F_{j+r-1}^{i-1}$ and hence $\mathrm{d}(F_{j+\omega}^{i-1}) \subseteq \mathrm{d}(F_{j+r-1}^{i-1})$.
    Thus $v \in F_j^i \cap \mathrm{d}(F_{j+r-1}^{i-1}) = B_j^i(r)$, which means the class of $v$ is zero in $E_j^i(r)$.
    Therefore $\mathrm{d}_r = 0$ for all $r \ge \omega+1$, and the spectral sequence collapses at the $(\omega+1)$-page; i.e., $E^{i}_{j}(\omega+k)\simeq E^{i}_{j}(\omega+1)$ for all $k\ge 1$.

    By Proposition~\ref{prop:Boardman's_convergence}, we obtain that $\frac{\overline{Z}^{i}_{j}(\infty)}{\overline{B}^{i}_{j}(\infty)}\simeq \mathrm{gr}_{\bullet}^{F}H^{i}(F)$, take $r\gg 0$, we consider the following map:

    \begin{center}
        \begin{tikzcd}
            {\frac{\overline{Z}_j^i(r)}{\overline{B}_j^i(r)}} & {\frac{\overline{Z}_j^i(\infty)}{\overline{B}_j^i(r)}} & {\frac{\overline{Z}_j^i(\infty)}{\overline{B}_j^i(\infty)}} \\
            & {\frac{\overline{Z}_j^{i+1}(r+k)}{\overline{B}_j^{i+1}(r+k)}}
            \arrow["\simeq", from=1-1, to=2-2, sloped, dash]
            \arrow["\Phi"', hook', from=1-2, to=1-1]
            \arrow["\Psi", two heads, from=1-2, to=1-3]
            \arrow[from=1-2, to=2-2]
        \end{tikzcd}
    \end{center}

    1. Since $\mathrm{d}_k = 0$ for all $k \ge r$, we have $\frac{\overline{Z}_j^i(k)}{\overline{B}_j^i(k)} = \frac{\overline{Z}_j^i(k+1)}{\overline{B}_j^i(k+1)}$ for $k \ge r$.
    Since we have tower
    \[
        0 \subseteq \overline{B}_j^i(0) \subseteq \overline{B}_j^i(1) \subseteq \cdots \subseteq \overline{B}_j^i(r) \subseteq \cdots \subseteq \overline{B}_j^i(\infty)
        \subseteq \overline{Z}_j^i(\infty) \subseteq \cdots \subseteq \overline{Z}_j^i(r) \subseteq \overline{Z}_j^i(r-1) \subseteq \cdots \subseteq \overline{Z}_j^i(0)
    \]
    This implies the filtration on $Z$ and $B$ stabilize:
    $$ \overline{Z}_j^i(r) = \overline{Z}_j^i(r+1) = \cdots = \overline{Z}_j^i(\infty) $$
    and
    $$ \overline{B}_j^i(r) = \overline{B}_j^i(r+1) = \cdots = \overline{B}_j^i(\infty). $$
    Thus the map $\Phi$ (induced by inclusion $\overline{Z}_j^i(\infty) \subseteq \overline{Z}_j^i(r)$) is an isomorphism because the inclusion is actually an equality.

    2. Similarly, the map $\Psi$ (induced by projection modulo $\overline{B}_j^i(\infty)$ vs $\overline{B}_j^i(r)$) is an isomorphism because the denominators are equal.

    Therefore, $E_j^i(k)\simeq\frac{\overline{Z}^{i}_{j}(\infty)}{\overline{B}^{i}_{j}(\infty)} \simeq \mathrm{gr}_j H^i(F) $ for all $j$.
\end{proof}

\begin{remark}\label{rmk:classical_convergence_theorem}
    We introduce the \textbf{Classical Convergence Theorem} \cite[Classical~Convergence~Theorem~5.5.1]{Weibel1994HA} for a complex with \textbf{bounded below filtration} in the sense of convergence introduced in \cite[Convergence~5.2.11]{Weibel1994HA}.

    We will use this theorem in accompany by Lemma~\ref{lemma 7.8}.
\end{remark}

\begin{theorem}\label{tehorem 7.9}
    If $F^\bullet$ is $\omega$-regular and convergent in \cite[Convergence~5.2.11]{Weibel1994HA}'s sense, then $\mathrm{gr}^F_{\bullet} H^{i}(F^{\bullet})$ is a subquotient of $H^{i} \mathrm{gr}^{F}_{\bullet }(F^{\bullet})$.
\end{theorem}
\begin{proof}
    By Lemma \ref{lemma 7.4} + Lemma \ref{lemma 7.8}.
\end{proof}

\section{More on filtered rings }\label{sect:more_on_filtered_rings}

Recall that we have defined a $\mathbb{Z}$-filtration of a ring in Definition \ref{def:z_filtration}.

Let $A$ be a ring and $\Gamma_{\bullet} = \Gamma_{\bullet}A$ a $\mathbb{Z}$-filtration on $A$, recall that $\Gamma_{\bullet}A$ is called a $\textbf{noetherian } \mathbb{Z}\textbf{-filtration}$ if additionally $\mathrm{gr}^{\Gamma}_\bullet A$ is also noetherian.

From now on, we assume that $A$ is a $\C$-algebra and has a noetherian $\mathbb{Z}$-filtration $\Gamma_{\bullet}$, so it's rees algebra $R_{\bullet}^{\Gamma}A$ is a $\C[u]$ algebra.

Let $M$ be a finitely-generated $A$-module, and $\Omega_{\bullet}M$ a good filtration on $M$ s.t. $\mathrm{gr}^{\Omega} M \neq 0$.
Since $\mathrm{gr}^{\Gamma} A$ is noetherian, $R^{\Omega} M$ has a graded free resolution:
$$ \dots\rightarrow P_{\bullet}^{-1}\rightarrow P^{0}_{\bullet}\rightarrow R^{\Omega} _{\bullet}M\rightarrow0$$
s.t. each $P^{i}_{\bullet}$ is a free graded $R_{\bullet}^{\Gamma}A$-module of finite rank.
So in derived category, $\mathrm{RHom}_{R_{\bullet}A}(R_{\bullet}M, R_{\bullet}A)$ is the complex \cite[\href{https://stacks.math.columbia.edu/tag/0A8H}{Tag 0A8H}, \href{https://stacks.math.columbia.edu/tag/0A66}{Tag 0A66}]{stacks-project} (well-defined up to quasi-isomorphism in $\mathbf{D}(R_{\bullet}A)$):
$$\mathrm{Hom}_{R_{\bullet}A}(P^{\bullet}_{\bullet}, R_{\bullet}A).$$
Tensoring with $\mathbb{C}[u]/(u-1)$ over $\mathbb{C}[u]$, we get a free resolution:
$$ \cdots \to F^{-1} \to F^0 \to M \to 0$$
with $F^i = P^i_{\bullet} \otimes_{\mathbb{C}[u]} \mathbb{C}[u]/(u-1)$.
Define the filtration on $F^i_{\bullet}$ by:$$F_j^i := \mathrm{Im}(P_j^i \longrightarrow F^i)$$
This filtered complex:
$$\cdots \longrightarrow F_{\bullet}^{-1} \longrightarrow F_{\bullet}^0 \longrightarrow 0$$
is called a $\textbf{filtered free resolution}$ of $M$, which is $\C[u]/(u-1)\otimes^{L}_{\C[u]} R^{\Omega}M = A\otimes^{L}_{R_{\bullet}^{\Gamma}A} R^{\Omega}M$ in derived category, \cite[\href{https://stacks.math.columbia.edu/tag/0FNB}{Tag 0FNB}, \href{https://stacks.math.columbia.edu/tag/064M}{Tag 064M}]{stacks-project}
so it is well-defined up to quasi-isomorphism in $\mathbf{D}(A)$.

\bigskip

Since $\mathrm{Hom}({-},A)$ functor is contravariant, we define a filtered complex $$ \check{F}^{\bullet} = \operatorname{Hom}_{A}(F^{\bullet}, A) \quad \mathrm{s.t.}\quad \check{F}^{i} = \operatorname{Hom}_{A}(F^{-i}, A) $$
We consider $\mathrm{Hom}_{R_{\bullet}A}(P^{\bullet}_{\bullet},R_{\bullet}A)\otimes \mathbb{C}[u] /(u-1)$,
since $P^\bullet_{\bullet}$ is a free $R_{\bullet}A$-module,
we then have
$$\mathrm{Hom}_{R_{\bullet}A}(P^\bullet_{\bullet},R_{\bullet}A)\otimes \mathbb{C}[u]/(u-1)\simeq \mathrm{Hom}_{A}(F^{\bullet},A) \quad (\text{quasi-isomorphism in } \mathbf{D}(A)),$$
and filtration on $\check{F}^{\bullet}$ is defined as follows.

Recall that for two graded modules $M_{\bullet}, N_{\bullet}$ over a graded ring $R_{\bullet}$, the group of graded morphisms of degree $k$ is
(see \cite[\href{https://stacks.math.columbia.edu/tag/09MN}{Tag 09MN}]{stacks-project})
$$\mathrm{Hom}_{R_{\bullet}}(M_{\bullet}, N_{\bullet})_k := \{ f: M_{\bullet} \to N_{\bullet} \mid f(M_j) \subseteq N_{j+k}, \forall j \}.$$
Let $\mathcal{H}^i_{\bullet} := \mathrm{Hom}_{R_{\bullet}A}(P^{-i}_{\bullet}, R_{\bullet}A)$, which is a graded $R_{\bullet}A$-module. The isomorphism above induces a map from each graded piece:
$$ \Phi_k: \mathcal{H}^i_k \longrightarrow \mathrm{Hom}_A(F^{-i}, A). $$
We define the filtration on $\check{F}^i$ by the image of this map:
$$ \check{F}^i_k := \mathrm{Im}(\Phi_k). $$

\begin{proposition}\label{prop:filtration-on-Hom}
    The filtration defined above coincides with the standard induced filtration on $\mathrm{Hom}_A(F^{-i}, A)$, i.e.,
    $$ \check{F}^i_k = \{ f \in \mathrm{Hom}_A(F^{-i}, A) \mid f(F_j^{-i}) \subseteq F_{j+k}A, \forall j \}. $$
    (cf. \cite[Appendix~D, after Lemma~D.2.2]{hotta2008d})
\end{proposition}
\begin{proof}
    It suffices to check this for $P^{-i} = R_{\bullet}A(p)$, a shifted free module.
    Then $F^{-i} \cong A$ with filtration shifted by $p$, i.e., $F_j^{-i} = F_{j+p}A$.

    LHS: $\mathcal{H}^i_k = \mathrm{Hom}_{R_{\bullet}A}(R_{\bullet}A(p), R_{\bullet}A)_k \cong R_{\bullet}A_{k-p}$.
    Its image in $\mathrm{Hom}_A(A, A) \cong A$ is exactly $F_{k-p}A$.

    RHS: $\{ f: A \to A \mid f(F_j^{-i}) \subseteq F_{j+k}A \} = \{ a \in A \mid a \cdot F_{j+p}A \subseteq F_{j+k}A, \forall j \}$.
    Since $A$ has a filtration, this condition is equivalent to $a \in F_{k-p}A$ (by definition of filtration on ring).

    Thus LHS = RHS. Since $P^{-i}$ is a direct sum of such terms, the result holds.
\end{proof}

\begin{lemma}\label{lem:checkF-regular}
    $\check{F}^{\bullet}$ is  $\omega$-regular for some $\omega \gg 0.$
\end{lemma}
\begin{proof}
    Take $N = \mathrm{d}(\check{F}^{k}) \subseteq \check{F}^{k+1}$, then $N_j = \mathrm{d}(\check{F}_{j}^{k})$ gives a good filtration on $N$.
    On the other hand, $N \cap \check{F}_{j}^{k+1}$ gives another good filtration. By Corollary \ref{cor:compare-good-filtrations}, there exists a uniform $\omega$ s.t. $$ \check{F}_{j}^{k+1} \cap \mathrm{d}(\check{F}^{k}) = N \cap \check{F}_{j}^{k+1} \subseteq N_{j+\omega} = \mathrm{d}(\check{F}_{j+\omega}^{k}), \quad \forall j. $$
\end{proof}

\begin{corollary}\label{corollary 7.11}
    $\mathrm{gr}_{\bullet}\mathrm{Ext}_A^j (M, A)$ is a subquotient of $\mathrm{Ext}_{\mathrm{gr}_{\bullet} A}^j (\mathrm{gr} _{\bullet}M, \mathrm{gr}_{\bullet} A)$ for all $j$.
\end{corollary}
\begin{proof}
    Applying Theorem \ref{tehorem 7.9} to the filtered complex $\check{F}^{\bullet}$.
    \begin{itemize}
        \item By Lemma \ref{lem:checkF-regular}, $\check{F}^{\bullet}$ is $\omega$-regular.
        \item Since the filtration on $M$ is a good filtration, it is bounded from below (see \cite[Lemma~D.2.3]{hotta2008d}). This implies the filtration on $\check{F}^\bullet$ is also bounded from below, so the spectral sequence converges (see Remark~\ref{rmk:classical_convergence_theorem}).
    \end{itemize}
    Thus $\mathrm{gr}_{\bullet}H^j(\check{F}^{\bullet})$ is a subquotient of $H^j(\mathrm{gr}_{\bullet}\check{F}^{\bullet})$.

    \medskip

    Now we identify the terms:
    \begin{enumerate}
        \item By definition, $H^j(\check{F}^{\bullet}) = \mathrm{Ext}_A^j(M, A)$. Thus $\mathrm{gr}_{\bullet}H^j(\check{F}^{\bullet}) = \mathrm{gr}_{\bullet}\mathrm{Ext}_A^j(M, A)$.
        \item Since $F^\bullet$ comes from a graded free resolution $P^\bullet$, $\mathrm{gr}_{\bullet} F^{\bullet} \cong P^\bullet / u P^\bullet$ is a free resolution of $\mathrm{gr}_{\bullet} M \cong R^\Omega M / u R^\Omega M$ over $\mathrm{gr}_{\bullet} A \cong R^\Gamma A / u R^\Gamma A$.
        \item Since $F^k$ are free modules, taking the associated graded commutes with $\mathrm{Hom}$, i.e.,
              \[
                  \mathrm{gr}_{\bullet}\check{F}^{\bullet} = \mathrm{gr}_{\bullet}\mathrm{Hom}_A(F^{\bullet}, A) \text{ and } \mathcal{G}^{\bullet} = \mathrm{Hom}_{\mathrm{gr}_{\bullet}A}(\mathrm{gr}_{\bullet}F^{\bullet}, \mathrm{gr}_{\bullet}A).
              \]
              We have a canonical map of graded modules:
              \[ \Psi: \mathrm{gr}_{\bullet}\mathrm{Hom}_A(F^{\bullet}, A) \longrightarrow \mathrm{Hom}_{\mathrm{gr}_{\bullet}A}(\mathrm{gr}_{\bullet}F^{\bullet}, \mathrm{gr}_{\bullet}A). \]
              It is defined by the principal symbol: for $f \in \check{F}^k_p = F_p\mathrm{Hom}_A(F^k, A)$, its image $\Psi(f \bmod \check{F}^k_{p-1})$ is the graded morphism:
              \[
                  \mathrm{gr}_{\bullet}F^k \ni (x \bmod F_{j-1}F^k) \longmapsto (f(x) \bmod F_{j+p-1}A) \in \mathrm{gr}_{\bullet}A.
              \]
              Since $F^k$ is a filtered free module (direct sum of shifted $A$), and for $A(s)$, the map $\Psi$ corresponds to the identity on $(\mathrm{gr}_{\bullet}A)(-s)$. Thus $\Psi$ is an isomorphism (or see \cite[Lemma~D.2.3]{hotta2008d}).
        \item Therefore, $H^j(\mathrm{gr}_{\bullet}\check{F}^{\bullet}) \cong \mathrm{Ext}_{\mathrm{gr}_{\bullet}A}^j(\mathrm{gr}_{\bullet}M, \mathrm{gr}_{\bullet}A).$
    \end{enumerate}
    The result follows.
\end{proof}

\begin{remark}[{Induced filtration on $\Ext$ (cf. \cite[Appendix~D, in the proof of Lemma~D.2.4]{hotta2008d})}]\label{rmk:induced-filtration-on-Ext}
    \mbox{}

    To make sense of the spectral sequence convergence, we define the filtration on $\mathrm{Ext}$ group.

    Using the filtration on $\check{F}^\bullet$ defined above, the \textbf{induced filtration} on
    \[
        \Ext_A^i(M,N)\cong H^i(\check{F}^\bullet)
    \]
    is defined by
    \[
        F_p\Ext_A^i(M,N)
        := F_pH^i(\check{F}^\bullet)
        := \mathrm{Im}\!\left(H^i(F_p\check{F}^\bullet)\longrightarrow H^i(\check{F}^\bullet)\right),
    \]
    where $F_p\check{F}^\bullet$ denotes the subcomplex with $(F_p\check{F}^\bullet)^i=F_p\check{F}^i$.
    Note that this definition coincides with Definition \ref{def:induced_filtration_cohomology}.

    Indeed, the map $H^i(F_pK^\bullet) \to H^i(K^\bullet)$ sends the class $[z]_p$ of a cycle $z \in Z_p^i(\infty) = \ker(d) \cap F_p K^i$ to its class $[z]$ in $H^i(K^\bullet)$.
    Thus the image is
    \[
        \mathrm{Im}\!\left(H^i(F_pK^\bullet)\longrightarrow H^i(K^\bullet)\right) = \left\{ z + \mathrm{d}K^{i-1} \in H^i(K^\bullet) \mid z \in Z_p^i(\infty) \right\} = \frac{Z_p^i(\infty) + \mathrm{d}K^{i-1}}{\mathrm{d}K^{i-1}},
    \]
    which matches the definition above.
\end{remark}


\section{Duality for $\mathcal{D}_{X}$-mod}

Let $X = \mathrm{Spec}\, R$ a smooth affine variety with algebraic coordinates $(x_1, \ldots, x_n)$ and $M$ a nonzero f.g. $\mathcal{D}_X$-module.

\begin{theorem}\label{theorem 7.12}
    We have:
    \begin{enumerate}
        \item $\dim \operatorname{Ch}(\operatorname{Ext}_{\mathcal{D}_X}^i(M, \mathcal{D}_X)) \leq 2n - i$.
        \item $\operatorname{Ext}_{\mathcal{D}_X}^i(M, \mathcal{D}_X) = 0$ for $i < 2n - \dim(\operatorname{Ch}(M))$ or $i > n$.
        \item $j(M) + \dim \operatorname{Ch}(M) = 2n$.
    \end{enumerate}
\end{theorem}
\begin{proof}
    \textbf{(1) and (2)}
    By Corollary~\ref{corollary 7.11}, and the fact that the support of a subquotient is contained in the support of the original module, we have:
    $$ \operatorname{Ch}(\operatorname{Ext}_{\mathcal{D}_X}^i(M, \mathcal{D}_X)) = \operatorname{Supp}(\operatorname{gr}_{\bullet}(\operatorname{Ext}_{\mathcal{D}_X}^i(M, \mathcal{D}_X))) \subseteq \Ch(\mathrm{Ext}^{j}_{\mathrm{gr}_{\bullet}\mathcal{D}_{X}}(\mathrm{gr}_{\bullet}M,\mathrm{gr}_{\bullet}\mathcal{D}_{X})). $$
    Then apply Theorem \ref{theorem 7.2} to $\Ch(\mathrm{Ext}^{j}_{\mathrm{gr}_{\bullet}\mathcal{D}_{X}}(\mathrm{gr}_{\bullet}M,\mathrm{gr}_{\bullet}\mathcal{D}_{X}))$, since $\gr_{\bullet}\mathcal{D}_{X}$ is commutative.

    \medskip

    \textbf{(3)}  We know that $j(\mathrm{gr}_{\bullet}M)+\dim(\operatorname{Ch}(M)) = \dim R[\xi_{1},\dots,\xi_{n}] = 2n$ by Theorem \ref{theorem 7.2}(2). Thus, it suffices to show that $j(M)=j(\mathrm{gr}_{\bullet}M)$.

    Since by Theorem \ref{tehorem 7.9}, $\mathrm{gr}_{\bullet}\mathrm{Ext}_{\mathcal{D}_{X}}^{j}(M,\mathcal{D}_{X})$ is a subquotient of $\mathrm{Ext}^{j}_{\mathrm{gr_{\bullet}\mathcal{D}_{X}}}(\mathrm{gr}_{\bullet}M,\mathrm{gr}_{\bullet}\mathcal{D}_{X})$.
    Thus we obtain the inequality:
    \[
        j(M) \ge j(\mathrm{gr}_{\bullet}M) = 2n - \dim(\operatorname{Ch}(M)).
    \]

    Now, assume $j(\mathrm{gr}_{\bullet}M)=j_{0}$. This means $\mathrm{Ext}^{j_{0}}_{\mathrm{gr_{\bullet}\mathcal{D}_{X}}}(\mathrm{gr}_{\bullet}M,\mathrm{gr}_{\bullet}\mathcal{D}_{X})\neq 0$.
    Consider the spectral sequence (of $\check{F}^\bullet$) converging to $\mathrm{Ext}_{\mathcal{D}_{X}}^{\bullet}(M,\mathcal{D}_{X})$ (resulting from the $\omega$-regularity of $\check{F}^\bullet$,
    see Lemma \ref{lemma 7.8}, Remark~\ref{rmk:classical_convergence_theorem} and Lemma \ref{lem:checkF-regular}).
    We have $E^{j}_{\bullet}(1) \simeq \mathrm{Ext}^{j}_{\mathrm{gr}_{\bullet}\mathcal{D}_{X}}(\mathrm{gr}_{\bullet}M,\mathrm{gr}_{\bullet}\mathcal{D}_{X})$ (the filtration on $\mathrm{Ext}$ mentioned in Rmk.~\ref{rmk:induced-filtration-on-Ext} is the same as the one in Def.~\ref{def:induced_filtration_cohomology}).
    \begin{itemize}
        \item For $j < j_0$, $E^{j}_{\bullet}(1) = 0$.
        \item Thus $E^{j}_{\bullet}(k+1) \simeq \ker(E^{i}(k)\rightarrow E^{i+1}(k))$ for $k \ge 1$.
    \end{itemize}

    The order filtration on $\mathrm{gr}_{\bullet}\mathcal{D}_{X}$ is a $\mathbb{Z}_{\ge 0}$-filtration. Since $\Omega_{\bullet}M$ is a good filtration, it is bounded below (i.e., $\Omega_{p}M=0$ for $p \ll 0$). Consequently, the induced filtration on the resolution $\check{F}^{\bullet}$ is also bounded below.

    Let $p_{0}$ be the minimal integer such that $\check{F}^{j_{0}}_{p_{0}} \neq 0$.
    Then:
    \[
        E^{j_{0}}_{p_{0}}(0) = \check{F}^{j_{0}}_{p_{0}} \neq 0.
    \]
    For $k \ge 1$, the differential $\mathrm{d}_{k}: E^{j_{0}}_{p_{0}}(k) \rightarrow E^{j_{0}+1}_{p_{0}-k}(k)$ must be the zero map (since $p_0-k < p_0$, the target term is 0).
    Therefore,
    \[
        E^{j_{0}}_{p_{0}}(\infty) = E^{j_{0}}_{p_{0}}(0) \neq 0.
    \]
    This implies $\mathrm{gr}_{p_{0}}\mathrm{Ext}_{\mathcal{D}_{X}}^{j_{0}}(M,\mathcal{D}_{X}) \neq 0$, so $\mathrm{Ext}_{\mathcal{D}_{X}}^{j_{0}}(M,\mathcal{D}_{X}) \neq 0$.

    Hence $j(M) \le j_{0} = j(\mathrm{gr}_{\bullet}M)$. Combining with the reverse inequality, we get $j(M) = j(\mathrm{gr}_{\bullet}M)$.
\end{proof}

\begin{proposition}\label{prop:DX-has-finite-global-dimension}
    The global dimension of $\mathrm{d}_X$ is also finite. Furthermore, the global dimension of $\mathcal{D}_X$ $=n$.
\end{proposition}
\begin{proof}
    First, we show $\operatorname{gl.dim}\mathcal{D}_X \ge n$. Take $M=\mathcal{O}_{X}$. Since $M$ is holonomic, $\dim(\operatorname{Ch}(M))=n$. Using Theorem \ref{theorem 7.12}(3), we have $j(M)=2n-n=n$, which implies $\operatorname{Ext}^n(M, \mathcal{D}_X) \neq 0$. Thus $\operatorname{pd}(M) \ge n$, so $\operatorname{gl.dim}\mathcal{D}_X \ge n$.

    Now we show $\operatorname{gl.dim}\mathcal{D}_X \le n$. It suffices to show that for any non-zero finitely generated $\mathcal{D}_X$-module $M$, $\operatorname{Ext}^i(M, \mathcal{D}_X) = 0$ for all $i > n$.
    Suppose $\operatorname{Ext}^i(M, \mathcal{D}_X) \neq 0$ for some $i > n$. By Theorem \ref{theorem 7.12}(1), we have:
    $$\dim \operatorname{Ch}(\operatorname{Ext}^i(M, \mathcal{D}_X)) \le 2n - i < 2n - n = n.$$
    However, by Bernstein's Inequality (Theorem \ref{thm:Bernstein-inequality}, which implies $\dim \operatorname{Ch}(M) \ge n$), the dimension of the characteristic variety must be at least $n$. This is a contradiction.
    Therefore, $\operatorname{Ext}^i(M, \mathcal{D}_X) = 0$ for all $i > n$, which means $\operatorname{pd}(M) \le n$. Consequently, $\operatorname{gl.dim}\mathcal{D}_X \le n$.
\end{proof}

\begin{theorem}
    $\mathcal{D}_{X}$ is Auslander regular.
\end{theorem}
\begin{proof}
    By Proposition \ref{prop:DX-has-finite-global-dimension}, $\operatorname{gl.dim}\mathcal{D}_X = n$, a finite number.

    If $N \subset \operatorname{Ext}^{i}(M, \mathcal{D}_X)$ is a nonzero submodule, then $\operatorname{Ch}(N)\subset\operatorname{Ch}(\operatorname{Ext}^{i}(M, \mathcal{D}_X))$.

    So $\dim(\operatorname{Ch}(N)) \leq \mathrm{dim}(\operatorname{Ch}(\operatorname{Ext}^{i}(M, \mathcal{D}_X)))\leq 2n - i$ by Theorem \ref{theorem 7.12}(1).
    Then $j(N) \geq i$ by Theorem \ref{theorem 7.12}(3).
\end{proof}

\begin{corollary}
    \quad
    \begin{enumerate}
        \item $M$ is holonomic iff $\operatorname{Ext}^i(M, \mathcal{D}_X) = 0$, $i \neq n$.
        \item If $M$ is holonomic, then $\operatorname{Ext}^n(M, \mathcal{D}_X)$ is also holonomic.
    \end{enumerate}
\end{corollary}
